\section{Fibonacci splitting of the Mersenne primes}
\label{sec:fibonacci}

The modular concentration of~\Cref{thm:concentration-main} records
\emph{where} the Mersenne primes fall in $\Ztwenty$: their residues lie
in $\{3,7,11\}$.  We now show \emph{how} they fall---how each prime
$\Mp$ splits in the field $\mathbb{Q}(\sqrt5)$ of the golden ratio.  The
answer is governed by a single classical congruence (Lucas, 1878) and,
composed with the concentration, refines the three admissible classes
into a split class and two inert classes.  This is the arithmetic
content behind the identification $3=M_{2}=|\{3,7,11\}|$ of
clause~(vii): the multiplier~$3$ is the number of classes the Mersenne
primes occupy, and the golden field determines the splitting type in
each.

\subsection{Fibonacci prerequisites}
\label{ssec:fib-prereq}

The Fibonacci sequence is $F_{0}=0$, $F_{1}=1$, $F_{n+2}=F_{n+1}+F_{n}$.
In any commutative ring containing an element~$\alpha$ with
$\alpha^{2}=\alpha+1$, induction on~$n$ gives the Binet-type identity
\begin{equation}\label{eq:binet}
  \alpha^{\,n+1} = F_{n+1}\,\alpha + F_{n}.
\end{equation}
The Galois conjugate of~$\golden$ over~$\mathbb{Q}$ is
$\overline{\golden}=1-\golden=-1/\golden$, with trace
$\golden+\overline{\golden}=1$ and norm
$\golden\,\overline{\golden}=-1$; both roots of $X^{2}-X-1$.  The
discriminant of $\mathbb{Q}(\sqrt5)/\mathbb{Q}$ is
$(\golden-\overline{\golden})^{2}=5$, so~$5$ is the unique ramified
prime.  For an odd prime $q\neq5$, the splitting of~$q$ in
$\mathbb{Q}(\sqrt5)$ is governed by the Legendre symbol $(q/5)$:
$(q/5)=+1$ (split) iff $q\bmod 20\in\{1,9,11,19\}$, and $(q/5)=-1$
(inert) iff $q\bmod 20\in\{3,7,13,17\}$.  Among the admissible Mersenne
classes~$\{3,7,11\}$, exactly one---the class~$11$---is split.
Lean: \lean{fib\_pow\_aux}, \lean{chi5\_values}.

\subsection{The Fibonacci splitting congruence}
\label{ssec:fib-congruence}

\begin{Theorem}[Fibonacci splitting; Lucas 1878]\label{thm:fib-main}
For every prime $q\neq2,5$,
\[
  F_{q}\;\equiv\;(q/5)\pmod{q},
\]
with value $+1$ in the split case and $-1$ in the inert case.
\emph{Lean:} \lean{fibonacci\_splitting\_general}.
\end{Theorem}

\begin{Lemma}[Split case]\label{lem:fib-split}
Let $q\neq2,5$ be prime with $(q/5)=+1$.  Then $X^{2}-X-1$ has a root
$\alpha\in\mathbb{F}_{q}$, and $F_{q}\equiv 1\pmod{q}$.
\emph{Lean:} \lean{fib\_split\_case}.
\end{Lemma}

\begin{proof}
Since $(q/5)=+1$, there is $\beta\in\mathbb{F}_{q}$ with $\beta^{2}=5$;
put $\alpha:=(1+\beta)/2$ (defined as $q\neq2$).  Then
$\alpha^{2}=(6+2\beta)/4=\alpha+1$.  By Fermat, $\alpha^{q}=\alpha$, and
by~\eqref{eq:binet}, $\alpha^{q}=F_{q}\alpha+F_{q-1}$; the same holds for
the conjugate root $\overline{\alpha}=1-\alpha$.  Subtracting the two
relations gives $F_{q}(\alpha-\overline{\alpha})\equiv
\alpha-\overline{\alpha}$; since
$(\alpha-\overline{\alpha})^{2}=5\not\equiv0$ (as $q\neq5$) the factor is
invertible, so $F_{q}\equiv1\pmod{q}$.
\end{proof}

\paragraph{Worked example ($q=11$, split).}
Since $11\equiv1\pmod5$, $(11/5)=+1$; in $\mathbb{F}_{11}$ the roots of
$X^{2}=X+1$ are $\alpha=8$, $\overline{\alpha}=4$ (as
$8^{2}=64\equiv9=8+1$).  And $F_{11}=89=11\cdot8+1$, so
$F_{11}\equiv1\pmod{11}$.

\begin{Lemma}[Inert case]\label{lem:fib-inert}
Let $q\neq2,5$ be prime with $(q/5)=-1$.  Then $X^{2}-X-1$ is
irreducible over~$\mathbb{F}_{q}$, the Frobenius $x\mapsto x^{q}$ of
$\mathbb{F}_{q^{2}}=\mathbb{F}_{q}[X]/(X^{2}-X-1)$ exchanges the two
roots $\alpha,\overline{\alpha}=1-\alpha$, and $F_{q}\equiv-1\pmod{q}$.
\emph{Lean:} \lean{fib\_inert\_case}.
\end{Lemma}

\begin{proof}
By quadratic reciprocity, $(5/q)=(q/5)=-1$, so~$5$ is a non-residue
and $X^{2}-X-1$ has no root in~$\mathbb{F}_{q}$; being quadratic it is
irreducible, and the quotient
$\mathbb{F}_{q^{2}}=\mathbb{F}_{q}[X]/(X^{2}-X-1)$ is a field of
order~$q^{2}$ with basis $\{1,\alpha\}$, where $\alpha$ is the image
of~$X$.  The Frobenius $\mathrm{Frob}_{q}\colon x\mapsto x^{q}$ is an
$\mathbb{F}_{q}$-automorphism (additivity is the freshman's dream in
characteristic~$q$; fixity of~$\mathbb{F}_{q}$ is Fermat), it permutes
the roots of $X^{2}-X-1$, and it is not the identity (else all of
$\mathbb{F}_{q^{2}}$ would lie in~$\mathbb{F}_{q}$); hence
$\alpha^{q}=\overline{\alpha}=1-\alpha$.  By~\eqref{eq:binet},
$\alpha^{q}=F_{q}\alpha+F_{q-1}$, so $(F_{q}+1)\alpha+(F_{q-1}-1)=0$ in
$\mathbb{F}_{q^{2}}$; as $\{1,\alpha\}$ is a basis, both coefficients
vanish, giving $F_{q}\equiv-1\pmod{q}$.
\end{proof}

\begin{Remark}[Lucas--Lehmer companion]\label{rem:fib-companion}
The inert case also yields $F_{q-1}\equiv1$, whence
$F_{q+1}=F_{q}+F_{q-1}\equiv0\pmod{q}$: the classical Fermat--Lucas
condition $q\mid F_{q+1}$ for inert primes, which underlies the index of
appearance in the Lucas--Lehmer primality test.
\end{Remark}

\paragraph{Worked example ($q=7$, inert).}
Since $7\equiv2\pmod5$ is a non-residue, $(7/5)=-1$ and $X^{2}-X-1$ is
irreducible over~$\mathbb{F}_{7}$.  Here $F_{7}=13\equiv-1\pmod{7}$, and
$F_{8}=21\equiv0\pmod7$, consistent with~\Cref{rem:fib-companion}.

\subsection{Splitting type of the Mersenne primes}
\label{ssec:fib-mersenne}

\begin{Theorem}[Residue class determines splitting type]\label{thm:residue-split}
Let~$p$ be a prime with $\Mp$ also prime, $\Mp\neq2,5$.  Then
$\Mp\bmod20\in\{3,7,11\}$, and
\[
  \Mp\equiv11\Rightarrow F_{\Mp}\equiv+1\ (\text{split}),
  \qquad
  \Mp\bmod20\in\{3,7\}\Rightarrow F_{\Mp}\equiv-1\ (\text{inert}),
\]
all congruences modulo~$\Mp$.
\emph{Lean:} \lean{mersenne\_fibonacci\_type}.
\end{Theorem}

\begin{proof}
Since $\Mp=2^{p}-1$ is odd and $2^{p}-1=5$ has no solution, the hypothesis $\Mp\neq2,5$ is automatic for every Mersenne prime.  The concentration is~\Cref{thm:concentration-main} at~$p$.
By~\Cref{thm:fib-main}, $F_{\Mp}\equiv(\Mp/5)\pmod{\Mp}$; and
$(\Mp/5)$ is read from $\Mp\bmod5$: $11\equiv1$ gives $+1$, while
$3\equiv3$ and $7\equiv2$ are non-residues, giving $-1$.
\end{proof}

\begin{table}[ht]\centering\footnotesize
\begin{tabular}{@{}cccccc@{}}
\toprule
$p$ & $\Mp$ & $\Mp\bmod20$ & $(\Mp/5)$ & $F_{\Mp}\bmod\Mp$ & type \\
\midrule
$3$  & $7$          & $7$  & $-1$ & $-1$ & inert \\
$5$  & $31$         & $11$ & $+1$ & $+1$ & split \\
$7$  & $127$        & $7$  & $-1$ & $-1$ & inert \\
$13$ & $8191$       & $11$ & $+1$ & $+1$ & split \\
$17$ & $131071$     & $11$ & $+1$ & $+1$ & split \\
$19$ & $524287$     & $7$  & $-1$ & $-1$ & inert \\
$31$ & $2147483647$ & $7$  & $-1$ & $-1$ & inert \\
\bottomrule
\end{tabular}
\caption{Fibonacci splitting type of $\Mp$ for the first seven Mersenne
primes with $p\geq3$.}\label{tab:mersenne-fib}
\end{table}

Specialised to each of the 52 known Mersenne exponents, this classifies
every known Mersenne prime by its splitting type; the distribution
is~\Cref{tab:splitting}.

\begin{table}[ht]\centering\footnotesize
\begin{tabular}{@{}cccc@{}}
\toprule
$\Mp\bmod20$ & Type in $\mathbb{Q}(\sqrt5)$ & $F_{\Mp}\bmod\Mp$ & Count \\
\midrule
$3$  & inert & $-1$ & $1$  \\
$7$  & inert & $-1$ & $19$ \\
$11$ & split & $+1$ & $32$ \\
\midrule
\multicolumn{3}{@{}r}{\textbf{total}} & \textbf{52} \\
\bottomrule
\end{tabular}
\caption{Splitting type in $\mathbb{Q}(\sqrt5)$ of each of the 52 known
Mersenne primes, by residue class $\Mp\bmod20$.  Verified in Lean by
\texttt{native\_decide} for $p\leq31$ and derived from
\lean{mersenne\_concentration\_general} otherwise.}\label{tab:splitting}
\end{table}

\begin{sloppypar}\small
\noindent\emph{Explicit lists.}  The single exponent in class~$3$ is
$p=2$ (inert).  The 19 exponents with $\Mp\equiv7\pmod{20}$ (inert) are
$3$, $7$, $19$, $31$, $107$, $127$, $607$, $1279$, $2203$, $4423$, $86\,243$, $110\,503$, $216\,091$, $756\,839$, $1\,257\,787$, $20\,996\,011$, $24\,036\,583$, $25\,964\,951$, $37\,156\,667$.
The 32 exponents with $\Mp\equiv11\pmod{20}$ (split) are
$5$, $13$, $17$, $61$, $89$, $521$, $2281$, $3217$, $4253$, $9689$, $9941$, $11\,213$, $19\,937$, $21\,701$, $23\,209$, $44\,497$, $132\,049$, $859\,433$, $1\,398\,269$, $2\,976\,221$, $3\,021\,377$, $6\,972\,593$, $13\,466\,917$, $30\,402\,457$, $32\,582\,657$, $42\,643\,801$, $43\,112\,609$, $57\,885\,161$, $74\,207\,281$, $77\,232\,917$, $82\,589\,933$, $136\,279\,841$.
\end{sloppypar}

\begin{Remark}[Classification under equivalence]\label{rem:classification}
Together, \Cref{thm:concentration-main} and~\Cref{thm:residue-split}
realise a classification of the Mersenne primes within the
pentagonal--cyclotomic geometry of~$\golden$: the multiplier
$3=M_{2}=|\{3,7,11\}|$ counts the admissible classes, and the golden
field partitions them into one split class ($11$) and two inert classes
($3,7$).  The observed split/inert split---$32$ against $20$ over the
current 52---is empirical; under conjectural equidistribution one would
expect a ratio nearer $1:2$, and the deviation over a sample of size~$52$
is not conclusive.
\end{Remark}


\subsection{The splitting type from the exponent}
\label{ssec:fib-exponent}

The residue $\Mp\bmod20$ of~\Cref{thm:residue-split} is itself fixed by the
exponent modulo~$4$, which reduces the whole classification to the
arithmetic of~$p$ alone.

\begin{Corollary}[Splitting type from $p\bmod4$]\label{cor:split-from-p}
Let $\Mp$ be a Mersenne prime.  Its splitting type in $\mathbb{Q}(\sqrt5)$
is determined by $p\bmod4$:
\[
  p\equiv1\!\!\pmod4\ \Rightarrow\ \text{split},
  \qquad
  p\equiv3\!\!\pmod4\ \text{or}\ p=2\ \Rightarrow\ \text{inert}.
\]
  \emph{Lean:} \lean{mersenne\_mod5\_from\_p\_mod4}.
\end{Corollary}

\begin{proof}
The residues $2^{p}\bmod5$ have period~$4$, equal to $2,4,3,1$ for
$p\equiv1,2,3,0\pmod4$.  Hence $\Mp=2^{p}-1$ modulo~$5$ depends only on
$p\bmod4$, taking the values $1,3,2$ for $p\equiv1,2,3\pmod4$, so
$(\Mp/5)=+1,-1,-1$ respectively.  By~\Cref{thm:residue-split} this sign is
the splitting type.  (The case $p\equiv0\pmod4$ gives $5\mid\Mp$ and is
excluded, since $\Mp$ is a prime $\neq5$.)
\end{proof}

Thus, among the admissible classes, the split class~$11$ is exactly
$p\equiv1\pmod4$ and the inert classes $3,7$ are $p=2$ and $p\equiv3\pmod4$.
The finer modulus~$8$ does not separate the type any further: $p\bmod4$ is
the exact controlling residue of the exponent.

\begin{Corollary}[Fibonacci rank characterises the inert primes]\label{prop:rank-inert}
For a Mersenne prime $\Mp\neq5$,
\[
  \Mp\mid F_{\Mp+1}
  \;\iff\;
  \Mp\ \text{is inert in }\mathbb{Q}(\sqrt5)
  \;\iff\;
  p\equiv3\!\!\pmod4\ \text{or}\ p=2.
\]
  \emph{Lean:} \lean{mersenne\_rank\_iff\_inert}.
\end{Corollary}

\begin{proof}
If $\Mp$ is inert then $\Mp\mid F_{\Mp+1}$ by~\Cref{rem:fib-companion}.  If
$\Mp$ is split, the roots $\alpha,\overline{\alpha}$ of $X^{2}-X-1$ lie in
$\mathbb{F}_{\Mp}^{\times}$, so $\alpha^{\Mp-1}=\overline{\alpha}^{\Mp-1}=1$
by Fermat, giving
$F_{\Mp-1}=(\alpha^{\Mp-1}-\overline{\alpha}^{\Mp-1})/(\alpha-\overline{\alpha})\equiv0$;
with $F_{\Mp}\equiv1$ (\Cref{lem:fib-split}) this yields
$F_{\Mp+1}=F_{\Mp}+F_{\Mp-1}\equiv1\not\equiv0\pmod{\Mp}$.  The second
equivalence is~\Cref{cor:split-from-p}.
\end{proof}

The divisibility $\Mp\mid F_{\Mp+1}$---that is, the rank of apparition of
$\Mp$ in the Fibonacci sequence dividing $\Mp+1$---therefore singles out
\emph{exactly} the inert half of the Mersenne primes, the same
rank-of-apparition structure that underlies the Lucas--Lehmer test.

\begin{Corollary}[The joint presentation from the exponent]\label{cor:joint}
For a Mersenne prime $\Mp$, the exponent residue $p\bmod4$ fixes the
triple (residue class modulo~$20$, splitting type in $\mathbb{Q}(\sqrt5)$,
Fibonacci sign $F_{\Mp}\bmod\Mp$) at once:
\[
  \begin{aligned}
    p\equiv1\!\!\pmod4 &\mapsto(11,\ \mathrm{split},\ +1),\\
    p\equiv3\!\!\pmod4 &\mapsto(7,\ \mathrm{inert},\ -1),\\
    p=2 &\mapsto(3,\ \mathrm{inert},\ -1).
  \end{aligned}
\]
Over the 52 known Mersenne primes the triple takes exactly these three
values, one per admissible residue; given $p\bmod4$, no further datum
distinguishes the class, the type, and the sign.
\begin{proof}
The class follows from the modular concentration
(\Cref{thm:concentration-main}) together with $\Mp\bmod5$, which
$p\bmod4$ fixes (\Cref{cor:split-from-p}); the type from the same
residue; and the Fibonacci sign from quadratic reciprocity, since
$X^{2}-X-1$ has a root modulo $\Mp$ exactly when $5$ is a square there,
i.e.\ when $\Mp\equiv\pm1\pmod5$, whence \Cref{lem:fib-split} or
\Cref{lem:fib-inert} applies.
\end{proof}
\emph{Lean:} \lean{mersenne\_presentation\_from\_p\_mod4}.
\end{Corollary}


